\documentclass[11pt,a4paper]{article}

\usepackage[T1]{fontenc}
\usepackage[utf8]{inputenc}
\usepackage[margin=1in]{geometry}
\usepackage{lmodern}
\usepackage{microtype}
\usepackage{setspace}
\usepackage{parskip}
\usepackage{amsmath,amssymb,amsfonts,mathtools}
\usepackage{booktabs}
\usepackage{longtable}
\usepackage{array}
\usepackage{tabularx}
\usepackage{xcolor}
\usepackage{hyperref}
\usepackage{listings}
\usepackage{enumitem}

\hypersetup{
  colorlinks=true,
  linkcolor=blue!50!black,
  urlcolor=blue!50!black,
  citecolor=blue!50!black
}

\lstdefinestyle{cfgstyle}{
  basicstyle=\ttfamily\small,
  breaklines=true,
  columns=fullflexible,
  keepspaces=true,
  frame=single,
  framerule=0.3pt,
  backgroundcolor=\color{black!2}
}

\setstretch{1.08}

\newcommand{\R}{\mathbb{R}}
\newcommand{\vect}[1]{\boldsymbol{#1}}
\newcommand{\mat}[1]{\mathbf{#1}}
\newcommand{\trans}{^{\mathsf T}}
\newcommand{\code}[1]{\texttt{#1}}

\title{Simulation and Control Integration Report\\
quad\_mini Quadruped in the OCS2--ROS2 Workspace}
\author{}
\date{\today}

\begin{document}
\maketitle
\tableofcontents
\newpage

\section{Introduction}

This report documents the \code{quad\_mini} control and simulation stack in the workspace rooted at
\code{/home/kaan/quad\_ocs2\_ws}. The objective is to present a coherent technical study of the
simulation and control architecture used for the \code{quad\_mini} platform, with emphasis on the
mathematical structure of the controller, the role of the configuration files, and the consistency
between the optimization layer and the MuJoCo model. Rather than limiting the discussion to a
parameter inventory, the report explains how the equations, constraints, solver settings, and
actuator definitions interact in the closed loop.

The stack under study is a multi-package ROS 2 control system composed of:
\begin{itemize}[leftmargin=*]
  \item \code{motion\_control}: centroidal-model MPC, whole-body control, constraints, and state handling,
  \item \code{mujoco\_simulator}: MuJoCo dynamics, rendering, contact detection, and actuator-side control application,
  \item \code{user\_command}: goal and gait command publication,
  \item \code{launch\_simulation}: launch-time orchestration,
  \item \code{legged\_msgs}: custom messages and services connecting the system.
\end{itemize}

The specific quad\_mini inputs used in this report are:
\begin{itemize}[leftmargin=*]
  \item \code{legged\_control/user\_command/config/quad\_mini/reference.info},
  \item \code{legged\_control/user\_command/config/quad\_mini/task.info},
  \item \code{legged\_control/user\_command/config/quad\_mini/simulation.info},
  \item \code{legged\_control/mujoco\_simulator/models/quad\_mini/urdf/robot.xml},
  \item \code{legged\_control/mujoco\_simulator/models/quad\_mini/urdf/robot\_rough.xml},
  \item \code{legged\_control/mujoco\_simulator/models/b1\_custom/urdf/robot.xml} for scale comparison.
\end{itemize}

The central engineering objective is to explain how the \code{quad\_mini} configuration turns a generic OCS2-based legged
control framework into a concrete closed-loop quadruped simulator. In particular, the report explains:
\begin{enumerate}[leftmargin=*]
  \item the system architecture and data path,
  \item the state, input, mode, and reference definitions,
  \item the nonlinear MPC problem as it is actually instantiated here,
  \item the weighted whole-body control layer that maps centroidal decisions to joint commands,
  \item the MuJoCo-side actuator and sensing setup,
  \item the parameter consistency checks and tuning implications.
\end{enumerate}

The report is organized as follows. Section 2 summarizes the baseline system understanding obtained
from the source tree and configuration files. Section 3 presents the main technical analysis in six
subsections. Section 4 gives conclusion and recommendations. Section 5 collects appendices and raw
reference snippets.

\section{Summary of Previous Work}

The initial review of the workspace and its supporting files established the following baseline facts.

\begin{itemize}[leftmargin=*]
  \item The workspace is not a minimal one-package controller. It vendors a substantial part of OCS2, Pinocchio, hpp-fcl, qpOASES,
        MuJoCo assets, and several robot examples, while the custom quadruped application logic lives in \code{legged\_control}.
  \item The active controller architecture is a centroidal nonlinear MPC plus whole-body control cascade:
        \[
        \text{simulator state or estimated state} \rightarrow \text{MPC} \rightarrow \text{WBC} \rightarrow \text{joint commands}.
        \]
  \item For the custom node \code{legged\_robot\_sqp\_mpc}, the instantiated high-level optimizer is \code{ocs2::SqpMpc}. Although
        \code{task.info} also contains DDP and IPM blocks, the inspected runtime path is SQP-based.
  \item The \code{quad\_mini} configuration is conservative relative to the larger \code{b1\_custom} model. It uses smaller joint torque
        limits, a lower base height, lower actuator-side PD gains, and a shorter prediction horizon.
  \item The available material already provided useful parameter inventories and cross-file checks, but a more integrated explanation of
        the control formulation and execution path was still needed.
\end{itemize}

These observations form the starting point of the present analysis. The following section develops the
controller description in a more systematic way, linking the configuration files to the corresponding
dynamic model, optimization problem, execution layer, and simulator behavior.

\section{New Technical Work}

\subsection{Control Architecture and System-Level Integration}

The \code{quad\_mini} control stack is a four-layer closed-loop system:
\begin{enumerate}[leftmargin=*]
  \item \textbf{Command layer:} \code{user\_command} publishes reference trajectories and gait updates.
  \item \textbf{Optimization layer:} \code{motion\_control} constructs and solves the switched nonlinear optimal control problem.
  \item \textbf{Execution layer:} the whole-body controller computes joint torque, position, and velocity commands.
  \item \textbf{Plant layer:} \code{mujoco\_simulator} integrates the dynamics and republishes state or sensor feedback.
\end{enumerate}

The closed-loop information flow is:
\[
\text{target command} \rightarrow \text{reference manager} \rightarrow \text{SQP-MPC policy}
\rightarrow \text{weighted WBC} \rightarrow \text{joint-level control law} \rightarrow \text{MuJoCo plant}.
\]

On the return path:
\[
\text{MuJoCo state/sensor data} \rightarrow \text{SystemObservation update} \rightarrow \text{MRT interface} \rightarrow \text{policy evaluation}.
\]

At runtime, the main ROS 2 interfaces relevant to the \code{quad\_mini} loop are shown in Table~\ref{tab:rosio}.

\begin{table}[ht]
\centering
\small
\begin{tabularx}{\textwidth}{>{\ttfamily}l >{\ttfamily}l >{\raggedright\arraybackslash}X}
\toprule
\textbf{Name} & \textbf{Type} & \textbf{Purpose} \\
\midrule
joint\_control\_data & topic & Whole-body controller output consumed by MuJoCo \\
simulator\_state\_data & topic & Ground-truth simulator state feedback \\
simulator\_sensor\_data & topic & Sensorized IMU/joint/contact feedback path \\
legged\_robot\_mpc\_observation & topic & Current observation published by motion\_control \\
legged\_robot\_mpc\_target & topic & Target trajectories published by user\_command \\
legged\_robot\_mpc\_mode\_schedule & topic & Runtime gait template updates \\
start\_control & service & Initial handshake that returns the simulator state \\
\bottomrule
\end{tabularx}
\caption{Core ROS 2 interfaces for the closed loop.}
\label{tab:rosio}
\end{table}

The simulator-side timing parameters from \code{simulation.info} are:
\begin{equation}
\Delta t_{\mathrm{sim}} = 0.001~\mathrm{s}, \qquad
f_{\mathrm{MPC}} = 40~\mathrm{Hz}, \qquad
f_{\mathrm{WBC}} = 200~\mathrm{Hz}, \qquad
f_{\mathrm{render}} = 60~\mathrm{Hz}.
\end{equation}
Thus the WBC runs exactly five times for every MPC update:
\begin{equation}
\frac{f_{\mathrm{WBC}}}{f_{\mathrm{MPC}}} = \frac{200}{40} = 5.
\end{equation}
This is cleaner than the B1 setup, where the same ratio is not an integer.

The launch process resolves robot-specific file paths through the \code{robot\_type} argument and starts the nodes in the order:
\begin{enumerate}[leftmargin=*]
  \item MuJoCo simulator,
  \item user command node,
  \item after a delay, the motion-control node.
\end{enumerate}

The initial control transition is not implicit. The controller explicitly calls the \code{start\_control} service, receives the current
simulator state, constructs the first \code{SystemObservation}, and only then requests the first MPC policy. This makes the initial
condition mathematically well defined:
\begin{equation}
\vect{x}(t_0) = \mathrm{CentroidalStateFromSimulator}(q(t_0),\dot{q}(t_0)).
\end{equation}

\subsection{Reference, Contact Modes, and Trajectory Generation}

The \code{quad\_mini} configuration uses a 24-dimensional OCS2 state and a 24-dimensional input. The state is
\begin{equation}
\vect{x} =
\begin{bmatrix}
\vect{v}_{\mathrm{com}} \\
\vect{k}/m \\
\vect{p}_b \\
\vect{\theta}_{zyx} \\
\vect{q}_j
\end{bmatrix}
\in \R^{24},
\end{equation}
where:
\begin{itemize}[leftmargin=*]
  \item $\vect{v}_{\mathrm{com}} \in \R^3$ is the normalized linear centroidal momentum term,
  \item $\vect{k}/m \in \R^3$ is the normalized angular centroidal momentum term,
  \item $\vect{p}_b \in \R^3$ is the base position,
  \item $\vect{\theta}_{zyx} \in \R^3$ is the base orientation in ZYX Euler coordinates,
  \item $\vect{q}_j \in \R^{12}$ is the vector of actuated joint angles.
\end{itemize}

The input is
\begin{equation}
\vect{u} =
\begin{bmatrix}
\vect{f}_{LF} \\
\vect{f}_{RF} \\
\vect{f}_{LH} \\
\vect{f}_{RH} \\
\dot{\vect{q}}_j
\end{bmatrix}
\in \R^{24},
\end{equation}
with four 3D contact force vectors and twelve joint velocities.

The contact mode is encoded by four stance flags:
\begin{equation}
\sigma = 8c_{LF} + 4c_{RF} + 2c_{LH} + c_{RH}, \qquad c_{\ell}\in\{0,1\},
\end{equation}
which yields a mode number in $\{0,\ldots,15\}$.

\paragraph{Reference posture.}
The reference file defines:
\begin{equation}
v_{\mathrm{disp}} = 0.50, \qquad
\omega_{\mathrm{yaw}} = 0.30, \qquad
h_{\mathrm{com}} = 0.19.
\end{equation}
The default joint posture is the symmetric crouched stance
\begin{equation}
\vect{q}_{j,\mathrm{def}} =
\begin{bmatrix}
0.00 & 0.17 & -0.92 & 0.00 & 0.17 & -0.92 &
0.00 & 0.17 & -0.92 & 0.00 & 0.17 & -0.92
\end{bmatrix}\trans.
\end{equation}

\paragraph{Initial mode schedule.}
The default startup schedule is purely stance:
\begin{equation}
\sigma(t) = \text{STANCE}, \qquad t \in [0,0.5]\ \text{at initialization},
\end{equation}
and the default template is also a single-mode stance template over one second. Therefore, the configuration is not starting from a
cyclic trot or pace gait. It starts from a statically safe support condition.

\paragraph{Goal command generation.}
The user command node converts
\[
\text{\code{goal: }}\Delta x\ \Delta y\ \Delta z\ \Delta \psi_{\deg}
\]
into a two-point target trajectory. Let the current base pose extracted from the latest observation be
\begin{equation}
\vect{p}_{b,0} = \begin{bmatrix} x_0 & y_0 & z_0 \end{bmatrix}\trans, \qquad
\vect{\theta}_{0} = \begin{bmatrix} \psi_0 & \theta_0 & \phi_0 \end{bmatrix}\trans.
\end{equation}
Then the target pose used by the controller is
\begin{equation}
\vect{p}_{b,f} =
\begin{bmatrix}
x_0 + \Delta x \\
y_0 + \Delta y \\
h_{\mathrm{com}} + \Delta z
\end{bmatrix},
\qquad
\vect{\theta}_{f} =
\begin{bmatrix}
\psi_0 + \Delta \psi \pi / 180 \\
\theta_0 \\
\phi_0
\end{bmatrix}.
\end{equation}

The arrival time is estimated from the command magnitudes:
\begin{equation}
t_f = t_0 +
\max\left(
\frac{\sqrt{(\Delta x)^2 + (\Delta y)^2}}{v_{\mathrm{disp}}},
\frac{|\Delta \psi|}{\omega_{\mathrm{yaw}}}
\right).
\end{equation}
So even before discussing the optimizer, the reference-generation layer already imposes a velocity-limited motion envelope.

\subsection{NMPC Formulation and Solver Design}

The active optimal control layer is a switched nonlinear MPC built around the centroidal model. In abstract form, the finite-horizon
problem solved at each update is:
\begin{align}
\min_{\vect{x}(\cdot),\vect{u}(\cdot)} \quad &
\int_{t_0}^{t_0+T}
\left[
\frac{1}{2}\left(\vect{x}-\vect{x}^{\star}\right)\trans \mat{Q}\left(\vect{x}-\vect{x}^{\star}\right) +
\frac{1}{2}\left(\vect{u}-\vect{u}^{\star}\right)\trans \mat{R}\left(\vect{u}-\vect{u}^{\star}\right)
\right]dt \\
\text{subject to} \quad &
\dot{\vect{x}} = \vect{f}(\vect{x},\vect{u},\sigma), \\
& \vect{g}(\vect{x},\vect{u},\sigma,t) = \vect{0}, \\
& \vect{h}(\vect{x},\vect{u},\sigma,t) \ge 0.
\end{align}

For \code{quad\_mini}, the configuration explicitly selects:
\begin{equation}
\text{\code{centroidalModelType = 0}},
\end{equation}
which means full centroidal dynamics rather than the single rigid body approximation.

\paragraph{Centroidal dynamics interpretation.}
At the modeling level, the flow map can be viewed as
\begin{equation}
\dot{\vect{x}} = \vect{f}(\vect{x},\vect{u})
\end{equation}
with centroidal momentum dynamics driven by contact forces:
\begin{equation}
\dot{\vect{h}} =
\sum_{i\in \mathcal{C}(\sigma)}
\begin{bmatrix}
\vect{f}_i \\
(\vect{r}_i - \vect{c}) \times \vect{f}_i
\end{bmatrix}
+
\begin{bmatrix}
m\vect{g} \\
\vect{0}
\end{bmatrix}.
\end{equation}
In the actual implementation, the code delegates the precise mapping and its derivatives to the OCS2 Pinocchio-based centroidal model.

\paragraph{Reference input.}
The nominal input is not simply zero. It is a weight-compensating contact-force distribution:
\begin{equation}
\vect{u}^{\star}(t) =
\begin{bmatrix}
\vect{f}_{LF}^{\star} \\
\vect{f}_{RF}^{\star} \\
\vect{f}_{LH}^{\star} \\
\vect{f}_{RH}^{\star} \\
\vect{0}
\end{bmatrix},
\qquad
\vect{f}_i^{\star} =
\begin{cases}
\dfrac{mg}{n_c}\vect{e}_z, & i \in \mathcal{C}(\sigma),\\[1mm]
\vect{0}, & i \notin \mathcal{C}(\sigma),
\end{cases}
\end{equation}
where $n_c$ is the number of stance feet.

\paragraph{State cost matrix.}
The diagonal entries of the \code{quad\_mini} state weight matrix are
\begin{equation}
\mathrm{diag}(\mat{Q}) =
\begin{bmatrix}
100,\ 100,\ 200,\ 100,\ 500,\ 1500,\ 500,\ 500,\ 3000,\ 100,\ 200,\ 2000, \nonumber \\
3,\ 3,\ 3,\ 3,\ 3,\ 3,\ 3,\ 3,\ 3,\ 3,\ 3,\ 3
\end{bmatrix}.
\end{equation}
This means the controller strongly penalizes:
\begin{itemize}[leftmargin=*]
  \item base height error through $Q_{zz}=3000$,
  \item roll-related body deviation through the large orientation weight,
  \item angular momentum deviation about the most sensitive axis.
\end{itemize}
In engineering terms, it is a stability-first posture-keeping objective.

\paragraph{Input cost matrix.}
The task file defines
\begin{equation}
\mathrm{diag}(\mat{R}_{\mathrm{task}}) =
0.70\times 10^{-3}
\begin{bmatrix}
\mathbf{1}_{12} \\
2000\mathbf{1}_{12}
\end{bmatrix},
\end{equation}
but the actual controller further maps the foot-velocity penalty into joint-velocity coordinates through the base-to-foot Jacobian:
\begin{equation}
\mat{R} =
\begin{bmatrix}
\mat{R}_f & \mat{0} \\
\mat{0} & \mat{J}_{bf}\trans \mat{R}_v \mat{J}_{bf}
\end{bmatrix}.
\end{equation}
Therefore, the effective joint-velocity regularization is configuration dependent even though the task file only shows a diagonal matrix.

\paragraph{Equality constraints.}
The controller enforces:
\begin{itemize}[leftmargin=*]
  \item zero contact force for swing feet,
  \item zero Cartesian foot velocity for stance feet,
  \item a swing-foot normal-velocity tracking condition based on the precomputed vertical spline.
\end{itemize}

For a stance foot $i$:
\begin{equation}
\vect{v}_{ee,i} = \vect{0}.
\end{equation}
For a swing foot $i$, the precomputation block effectively enforces:
\begin{equation}
\vect{e}_z\trans \vect{v}_{ee,i} - \dot{z}_i^{\star}(t) = 0,
\end{equation}
because \code{positionErrorGain = 0.0} in the current \code{quad\_mini} model settings.

\paragraph{Inequality constraints.}
The principal contact inequality is the soft friction-cone constraint:
\begin{equation}
h_i(\vect{f}_i) =
\mu F_{z,i} - \sqrt{F_{x,i}^2 + F_{y,i}^2 + \varepsilon} \ge 0,
\end{equation}
with
\begin{equation}
\mu = 0.8, \qquad
\mu_{\mathrm{barrier}} = 0.1, \qquad
\delta_{\mathrm{barrier}} = 5.0.
\end{equation}
The high friction coefficient compared to the B1 file means the \code{quad\_mini} configuration assumes fairly strong foot-ground traction.

\paragraph{Solver budget.}
The three solver blocks present in \code{task.info} are:
\begin{itemize}[leftmargin=*]
  \item SLQ/DDP with one iteration,
  \item SQP with one iteration and $dt=0.015$,
  \item IPM with one iteration and $dt=0.015$.
\end{itemize}
However, the active node instantiates SQP-MPC. Therefore, the runtime-relevant high-level numbers are:
\begin{equation}
T = 0.40~\mathrm{s}, \qquad
\Delta t = 0.015~\mathrm{s}, \qquad
N_{\mathrm{SQP}} = 1.
\end{equation}
This is a fast, lightweight, warm-start-dependent MPC regime. It is appropriate for fast iteration and debugging, but it leaves very
little room for deep per-step convergence.

\subsection{Whole-Body Control and Actuator-Level Realization}

The whole-body controller maps the centroidal-level desired state and input into generalized accelerations, contact forces, and joint
torques. The active implementation path is \code{WeightedWbc}. Its decision vector is
\begin{equation}
\vect{y} =
\begin{bmatrix}
\ddot{\vect{q}} \\
\vect{f}_c \\
\vect{\tau}
\end{bmatrix}
\in \R^{42},
\end{equation}
because the robot has:
\begin{equation}
18\ \text{generalized accelerations} +
12\ \text{contact-force variables} +
12\ \text{joint torques} = 42.
\end{equation}

\paragraph{Floating-base dynamics task.}
The rigid-body dynamics equality is
\begin{equation}
\mat{M}(\vect{q})\ddot{\vect{q}} + \vect{n}(\vect{q},\dot{\vect{q}})
= \mat{J}_c(\vect{q})\trans \vect{f}_c + \mat{S}\trans \vect{\tau}.
\end{equation}
This becomes one of the hard equality blocks of the WBC problem.

\paragraph{Contact acceleration task.}
For stance legs, the controller imposes acceleration-level no-slip conditions:
\begin{equation}
\mat{J}_{c,i}\ddot{\vect{q}} + \dot{\mat{J}}_{c,i}\dot{\vect{q}} = \vect{0}.
\end{equation}

\paragraph{Swing-leg task.}
For swing feet, the desired foot acceleration is generated through a PD law:
\begin{equation}
\vect{a}_{sw,i}^{\star}
= K_p\left(\vect{p}_{ee,i}^{\star} - \vect{p}_{ee,i}\right)
+ K_d\left(\dot{\vect{p}}_{ee,i}^{\star} - \dot{\vect{p}}_{ee,i}\right),
\end{equation}
with
\begin{equation}
K_p = 150, \qquad K_d = 12.
\end{equation}
This is lower and softer than the B1 swing gains, which is consistent with the smaller and lighter quad\_mini platform.

\paragraph{Friction pyramid and torque limits.}
At the WBC layer, the friction coefficient is again taken as
\begin{equation}
\mu_{\mathrm{WBC}} = 0.8,
\end{equation}
and the actuator torque limits are
\begin{equation}
\tau_{\max} =
\begin{bmatrix}
6 & 12 & 12
\end{bmatrix}\trans
\end{equation}
for the HAA, HFE, and KFE joints of each leg.

\paragraph{Weighted least-squares structure.}
The controller solves a QP of the form
\begin{align}
\min_{\vect{y}} \quad &
\frac{1}{2}\vect{y}\trans \mat{H}\vect{y} + \vect{g}\trans\vect{y} \\
\text{subject to} \quad &
\vect{l}_A \le \mat{A}\vect{y} \le \vect{u}_A,
\end{align}
where
\begin{equation}
\mat{H} = \mat{A}_w\trans\mat{A}_w, \qquad
\vect{g} = -\mat{A}_w\trans \vect{b}_w.
\end{equation}
The task weights in \code{task.info} are:
\begin{equation}
w_{\mathrm{swing}} = 100, \qquad
w_{\mathrm{base}} = 1, \qquad
w_{\mathrm{force}} = 0.01.
\end{equation}
Because the implementation multiplies the task residuals before forming $\mat{A}_w\trans\mat{A}_w$, the effective quadratic influence
scales like $w^2$, not $w$. Thus the swing-leg task is given overwhelmingly dominant importance relative to contact-force tracking.

\paragraph{Actuator-side realization in MuJoCo.}
The simulator does not apply the WBC torque output alone. Instead, each joint command is turned into a local control law:
\begin{equation}
\tau_i =
\tau_i^{ff}
+ K_p\left(q_i^{des} - q_i\right)
+ K_d\left(\dot{q}_i^{des} - \dot{q}_i\right),
\end{equation}
with
\begin{equation}
K_p = 7.5, \qquad K_d = 0.30
\end{equation}
from \code{simulation.info}. This is much gentler than the B1 simulator PD setup, which again matches the smaller robot scale.

\subsection{MuJoCo Simulation Model, Sensing, and Environment Variants}

\paragraph{Flat-ground quad\_mini model.}
The main MuJoCo model \code{robot.xml} defines:
\begin{itemize}[leftmargin=*]
  \item a flat plane contact surface,
  \item one floating base,
  \item twelve actuated joints,
  \item spherical feet of radius $0.02~\mathrm{m}$,
  \item IMU-like, joint, and base-frame sensors,
  \item twelve torque-limited motors.
\end{itemize}

The body mass distribution is approximately:
\begin{itemize}[leftmargin=*]
  \item trunk: $6.0~\mathrm{kg}$,
  \item lidar payload: $1.0~\mathrm{kg}$,
  \item camera payload: $0.4~\mathrm{kg}$,
  \item per leg: $0.39 + 0.91 + 0.30 + 0.01 = 1.61~\mathrm{kg}$.
\end{itemize}
Hence the total model mass is approximately
\begin{equation}
m \approx 6.0 + 1.0 + 0.4 + 4(1.61) = 13.84~\mathrm{kg}.
\end{equation}

\paragraph{Joint limits and actuator bounds.}
The motor ranges in the quad\_mini XML exactly match the WBC torque-limit task:
\begin{align}
\tau_{\mathrm{HAA}} &\in [-6,6], \\
\tau_{\mathrm{HFE}} &\in [-12,12], \\
\tau_{\mathrm{KFE}} &\in [-12,12].
\end{align}
This is an important consistency check because the optimization-side torque bounds and the simulator-side motor limits must agree to
avoid hidden saturation behavior.

\paragraph{Sensors.}
The sensor block publishes:
\begin{itemize}[leftmargin=*]
  \item base quaternion,
  \item base angular velocity,
  \item base linear acceleration,
  \item all joint positions,
  \item all joint velocities,
  \item base pose, base angular velocity, and base linear velocity through frame sensors.
\end{itemize}

The simulator converts these raw MuJoCo sensor values into:
\begin{equation}
\text{SimulatorStateData} \quad \text{or} \quad \text{SimulatorSensorData},
\end{equation}
depending on whether the control loop uses direct state feedback or estimator-like sensor feedback.

\paragraph{Rough-terrain variant.}
The file \code{robot\_rough.xml} preserves the same robot morphology and actuator limits, but it replaces the flat floor with a height field:
\begin{itemize}[leftmargin=*]
  \item ground type changes from \code{plane} to \code{hfield},
  \item the terrain is driven by \code{hfield.png},
  \item the same contact and actuator structure is retained.
\end{itemize}
This means the rough model is not a new robot. It is the same control plant embedded in a more demanding contact environment.

\paragraph{Comparison with b1\_custom.}
The \code{b1\_custom} model is much larger:
\begin{itemize}[leftmargin=*]
  \item trunk mass alone is $30.246~\mathrm{kg}$,
  \item hip and thigh links are also substantially heavier,
  \item torque limits rise to approximately $\pm 91.0$, $\pm 93.33$, and $\pm 140$.
\end{itemize}
This comparison matters because it explains why the \code{quad\_mini} gains, horizon, and torque limits cannot be interpreted as generic
legged-control settings. They are robot-scale dependent.

\subsection{Validation, Consistency Audit, and Engineering Interpretation}

The most important new engineering contribution of this report is not merely listing parameters, but interpreting how they interact across
files and across control layers.

\paragraph{Positive cross-file matches.}
The following consistency checks pass directly:
\begin{enumerate}[leftmargin=*]
  \item \textbf{Default posture consistency:} \code{reference.info.defaultJointState} matches \code{task.info.initialState} joint entries.
  \item \textbf{Foot radius consistency:} the MuJoCo foot radius of $0.02~\mathrm{m}$ matches \code{kalmanFilter.footRadius = 0.02}.
  \item \textbf{Torque-limit consistency:} the XML motor \code{ctrlrange} values match the WBC \code{torqueLimitsTask}.
  \item \textbf{Frequency consistency:} the quad\_mini WBC/MPC ratio is an exact integer, which simplifies multi-rate execution.
\end{enumerate}

\paragraph{Inactive or dormant features.}
Some configured items are present but not currently active for the inspected quad\_mini setup:
\begin{itemize}[leftmargin=*]
  \item \code{stateEstimate = false}, so the Kalman filter parameters are configured but dormant in the current loop.
  \item the task file contains DDP and IPM blocks, but the active node uses SQP-MPC.
  \item the self-collision block is effectively empty because the pair lists are not populated.
\end{itemize}

\paragraph{Stability-oriented interpretation.}
The combined effect of the current settings is clearly conservative:
\begin{itemize}[leftmargin=*]
  \item strong penalty on base height and orientation,
  \item short horizon of $0.40~\mathrm{s}$,
  \item one-iteration SQP budget,
  \item strong swing-leg priority in WBC,
  \item large foot-velocity penalty in the input cost,
  \item flat-stance default initial mode schedule.
\end{itemize}
Mathematically, this means the optimizer is encouraged to remain close to a well-supported crouched base posture while minimizing fast foot
motion relative to the nominal Jacobian frame. The result is a setup well suited for early integration and debugging, but not yet tuned for
the most aggressive dynamic behaviors.

\paragraph{Main leverage points for future tuning.}
If the goal is to make quad\_mini faster or more agile, the most meaningful tuning directions are:
\begin{enumerate}[leftmargin=*]
  \item increase the SQP iteration budget above $1$,
  \item consider extending the MPC horizon beyond $0.40~\mathrm{s}$,
  \item rebalance the very strong foot-velocity penalty embedded in $\mat{R}$,
  \item retune $K_p$, $K_d$, swing height, and lift-off/touchdown velocities for more dynamic swing trajectories,
  \item enable the estimator path only when the sensor model is validated strongly enough to justify the additional complexity.
\end{enumerate}

\section{Conclusion and Recommendations}

The \code{quad\_mini} configuration is internally coherent and technically well aligned with the custom OCS2--ROS2 architecture in this
workspace. The controller uses a full centroidal switched-model formulation, SQP-based nonlinear MPC, a weighted whole-body QP, and a
MuJoCo-side actuator law that adds compliant joint tracking around the WBC command. The reference design, torque limits, foot geometry,
and simulator motor bounds are mutually consistent, which is a strong sign that the configuration was assembled carefully rather than by
isolated parameter edits.

From a control-design perspective, the current setup should be understood as \textbf{stability-first and integration-first}, not
\textbf{agility-maximized}. The short horizon, one-iteration SQP budget, high base-stabilization weights, and strong foot-motion
regularization all point in that direction. This is a sensible design choice for bringing up a small quadruped in simulation.

The most direct recommendations are:
\begin{itemize}[leftmargin=*]
  \item keep the present configuration for baseline validation and regression testing,
  \item create a second, more aggressive parameter set rather than overwriting the stable one,
  \item treat solver-budget increase and input-cost rebalancing as the first two levers for more dynamic motion,
  \item use the rough-terrain XML as the first robustness benchmark before transferring any more aggressive controller settings.
\end{itemize}

\section{Appendices}

\subsection{Appendix A: Key Parameter Summary}

\begin{longtable}{p{0.38\linewidth}p{0.18\linewidth}p{0.36\linewidth}}
\toprule
\textbf{Domain} & \textbf{Value} & \textbf{Interpretation} \\
\midrule
\endhead
Centroidal model type & 0 & Full centroidal dynamics \\
State feedback mode & false & Direct simulator state, no estimator in active loop \\
MPC horizon & 0.40 s & Short prediction horizon \\
SQP discretization step & 0.015 s & Main MPC shooting step \\
Simulation timestep & 0.001 s & MuJoCo integration step \\
MPC frequency & 40 Hz & High-level optimizer rate \\
WBC frequency & 200 Hz & Mid-level torque computation rate \\
Render frequency & 60 Hz & Visualization rate \\
Target linear speed & 0.50 & Goal command velocity scale \\
Target yaw speed & 0.30 & Goal command yaw-rate scale \\
COM height reference & 0.19 m & Desired base height \\
Swing height & 0.03 m & Vertical foot lift amplitude \\
Lift-off velocity & 0.30 & Initial swing vertical speed \\
Touchdown velocity & -0.05 & Final swing descent speed \\
MPC friction coefficient & 0.8 & Soft friction-cone model \\
WBC friction coefficient & 0.8 & Friction pyramid in QP \\
WBC weights & 100 / 1 / 0.01 & Swing leg / base acceleration / contact force \\
Simulator PD gains & 7.5 / 0.30 & Joint-level tracking gains \\
Torque limits & 6 / 12 / 12 & HAA / HFE / KFE \\
\bottomrule
\end{longtable}

\subsection{Appendix B: Selected Raw Snippets}

\subsubsection*{B.1 Reference command block}
\begin{lstlisting}[style=cfgstyle]
targetDisplacementVelocity          0.50;
targetRotationVelocity              0.30;

comHeight                           0.19
\end{lstlisting}

\subsubsection*{B.2 MPC block}
\begin{lstlisting}[style=cfgstyle]
mpc
{
  timeHorizon                     0.40
  solutionTimeWindow              -1
  coldStart                       false
}
\end{lstlisting}

\subsubsection*{B.3 Whole-body control block}
\begin{lstlisting}[style=cfgstyle]
torqueLimitsTask
{
   (0,0)  6.0
   (1,0)  12.0
   (2,0)  12.0
}

swingLegTask
{
    kp                   150
    kd                   12
}

weight
{
    swingLeg        100
    baseAccel       1
    contactForce    0.01
}
\end{lstlisting}

\subsubsection*{B.4 Simulator timing and PD gains}
\begin{lstlisting}[style=cfgstyle]
controller
{
    timestep  0.001
    mpc_control_frequency  40
    wbc_control_frequency  200
}

pid
{
    kp                   7.5
    kd                   0.30
}
\end{lstlisting}

\subsubsection*{B.5 Rough-terrain floor}
\begin{lstlisting}[style=cfgstyle]
<geom name="floor"
      type="hfield"
      hfield="hfield"
      pos="0 0 -0.05"
      material="matplane"
      friction="1 0.005 0.0001"
      condim="3"/>
\end{lstlisting}

\subsection{Appendix C: File Coverage}

The main files used to support the analysis in this document were:
\begin{itemize}[leftmargin=*]
  \item \code{quad\_mini\_detailed\_report.tex},
  \item \code{quad\_ocs2\_project\_report.tex},
  \item \code{motion\_control/src/LeggedRobotSqpMpcNode.cpp},
  \item \code{motion\_control/src/ros\_interfaces/MPC\_WBC\_ROS\_Interface.cpp},
  \item \code{motion\_control/src/legged\_wbc/WbcBase.cpp},
  \item \code{motion\_control/src/legged\_wbc/WeightedWbc.cpp},
  \item \code{motion\_control/src/legged\_interface/LeggedRobotInterface.cpp},
  \item \code{motion\_control/src/robot\_precomputation/LeggedRobotPreComputation.cpp},
  \item \code{mujoco\_simulator/src/MujocoSimulation.cpp},
  \item \code{user\_command/src/goal/TargetTrajectoriesKeyboardPublisher.cpp},
  \item \code{user\_command/config/quad\_mini/*.info},
  \item \code{mujoco\_simulator/models/quad\_mini/urdf/*.xml},
  \item \code{mujoco\_simulator/models/b1\_custom/urdf/robot.xml}.
\end{itemize}

\end{document}
